% ============================================================
\section*{Conclusiones}
% ============================================================

\begin{itemize}
    \item Cifrar una contraseña no equivale a protegerla: el algoritmo elegido determina la
          diferencia entre una barrera criptográfica real y una falsa sensación de seguridad.
          SHA-256, diseñado para ser veloz, se convierte en una herramienta del atacante cuando
          se aplica al almacenamiento de contraseñas, pues su velocidad extrema ---más de $10^9$
          hashes por segundo en hardware moderno--- lo hace inútil como mecanismo de protección.
          En el Sistema IMC, el simple ajuste del \textit{cost factor} de BCrypt de~10 a~12
          cuadruplicó la resistencia criptográfica ($R_c$) sin cambiar el algoritmo, evidenciando
          que incluso decisiones paramétricas ---no solo la elección del algoritmo--- determinan
          si un sistema es realmente seguro o solo aparentemente seguro.

    \item Una implementación criptográfica técnicamente sofisticada puede quedar completamente
          anulada si el secreto que la sostiene está expuesto en el propio código fuente, visible
          para cualquier persona con acceso al repositorio. El Sistema IMC cometía exactamente
          este error: un secreto JWT hardcodeado y credenciales de administrador embebidas en
          \texttt{AuthService.java} invalidaban cualquier protección que BCrypt pudiera ofrecer,
          pues un atacante con acceso al código ya no necesitaba romper el cifrado sino simplemente
          leerlo. Externalizar las llaves a variables de entorno con validación de entropía mínima
          al arrancar transforma la seguridad criptográfica de una intención en una garantía
          técnica: si el secreto es débil o inexistente, la aplicación no arranca.

    \item Proteger las contraseñas en la base de datos mientras los tokens de sesión viajan en
          texto claro sobre HTTP es equivalente a blindar la puerta principal de una casa y dejar
          abierta la ventana trasera: la protección parcial genera una falsa sensación de seguridad
          que puede ser más peligrosa que su ausencia total, pues induce confianza donde no la hay.
          La configuración de TLS~1.3 con cipher suites que garantizan \textit{Perfect Forward
          Secrecy} cierra este vector de ataque, asegurando que incluso el compromiso futuro de
          la clave privada del servidor no exponga las sesiones pasadas. El Sistema IMC opera
          ahora con una arquitectura criptográfica coherente donde las capas de protección ---en
          reposo, en tránsito y en la gestión de llaves--- se complementan en lugar de
          contradecirse.
\end{itemize}
